%!TEX root = ../template.tex
%%%%%%%%%%%%%%%%%%%%%%%%%%%%%%%%%%%%%%%%%%%%%%%%%%%%%%%%%%%%%%%%%%%%
%% conclusion.tex
%% NOVA thesis document file
%%
%% Conclusions and Future Work chapter
%%%%%%%%%%%%%%%%%%%%%%%%%%%%%%%%%%%%%%%%%%%%%%%%%%%%%%%%%%%%%%%%%%%%

\typeout{NT FILE conclusion.tex}

\chapter{Conclusions and Future Work}
\label{cha:Conclusions}

\section{Conclusions}

This thesis presented the design and implementation of a cloud-based \gls{IoT} event management platform for monitoring and surveillance, extended with an edge computing node for offline resilience. The platform addresses the challenges inherent to \gls{IoT} monitoring systems (scalability, real-time data processing, security, and reliable operation in environments with intermittent connectivity) through a microservice architecture built on event-driven communication patterns.

The architecture was guided by five design goals. First, the platform was designed to be generic: it supports heterogeneous devices with varying telemetry formats through a schema-agnostic ingestion model that automatically discovers and normalizes telemetry fields without requiring predefined schemas. Second, real-time processing was achieved through an event-driven pipeline in which telemetry flows from device ingestion through analytics to user notification with minimal latency, using Kafka as a durable event backbone and WebSocket connections for live dashboard delivery. Third, extensibility was built into the architecture through the Kafka consumer group model, which allows new processing capabilities to be added by subscribing to existing topics without modifying producers or existing consumers. Fourth, security was addressed at every layer through a two-step device authentication flow, RS256 asymmetric \gls{JWT} signing for distributed token validation, per-device \gls{ACL} enforcement on both cloud and edge brokers, and centralized identity management through a dedicated authentication service. Fifth, offline resilience was achieved through the edge node's write-first pattern, which guarantees zero data loss during connectivity interruptions by persisting all telemetry and alerts to local SQLite storage before attempting cloud delivery.

The platform's cloud infrastructure consists of six microservices (an \gls{API} gateway, an authentication service, an ingestion service, an analytics service, a notifications service, and a read service) communicating through a dual-broker architecture with \gls{MQTT} for device-facing communication and Kafka for internal inter-service messaging. A Protocol Adapter in the communication layer provides \gls{HTTP} and \gls{CoAP} ingestion endpoints alongside the \gls{MQTT} broker, enabling devices that cannot use \gls{MQTT} to publish telemetry through alternative protocols that converge on the same Kafka pipeline. The Protocol Adapter also supports binary file uploads, storing files in MinIO and emitting metadata events to Kafka as an extensibility point for future processing capabilities. A lightweight form of the \gls{CQRS} pattern separates write-optimized and read-optimized responsibilities into distinct services, ensuring that high-throughput telemetry ingestion does not interfere with user-facing query performance. The platform provides dual detection capabilities: a rule-based analytics engine supporting ten comparison operators across numeric, string, and boolean data types, and a statistical anomaly detection engine using Z-score analysis with configurable sliding windows to identify deviations from normal behavior without requiring predefined thresholds. A rule severity system classifies rules as \texttt{CRITICAL}, \texttt{HIGH}, \texttt{MEDIUM}, or \texttt{LOW}. The notifications service provides email-based alert delivery with user-configurable severity filtering, consuming alert events independently from the same Kafka topic used for WebSocket delivery. A web application built with Next.js 16 provides device management, real-time telemetry dashboards, alert monitoring, and rule configuration through a browser-based interface.

The edge node, deployed on a Raspberry Pi 3B, extends the platform to environments where cloud connectivity cannot be guaranteed. Running a standalone NestJS application with a local Mosquitto broker and SQLite database, the edge node provides local telemetry ingestion, \texttt{CRITICAL} rule evaluation, and autonomous operation during network outages. The use of RS256 asymmetric \gls{JWT} signing is a key enabler for this architecture, allowing the edge node to validate device tokens locally using only the public key without requiring network access to the cloud Authentication Service. The edge node forwards pre-normalized telemetry and alerts to the cloud through dedicated \gls{MQTT} topics and Kafka consumers, ensuring that edge-originated data is unified with cloud data in the same storage and delivery pipeline.

The implementation is a proof-of-concept that demonstrates the viability of the architecture rather than a production-ready system. Infrastructure runs on single-node deployments without replication or clustering, internal communication lacks \gls{TLS} encryption, and credentials are managed through environment variables rather than a dedicated secret management system. The edge deployment supports a single Raspberry Pi node with no fleet management capabilities. These limitations are acknowledged as areas for future work rather than architectural deficiencies; the service boundaries, communication patterns, and data model are designed to support production hardening without structural changes.

Despite these limitations, the platform demonstrates that a modular, event-driven architecture can effectively address the core challenges of \gls{IoT} monitoring while remaining extensible for future capabilities. The edge node, in particular, validates that offline-resilient \gls{IoT} monitoring can be achieved on constrained hardware through careful resource management and a write-first data persistence strategy.

\section{Future Work}
\label{sec:FutureWork}

This section identifies directions for future development that were outside the scope of this proof-of-concept but would strengthen the platform for production deployment and broader adoption.

The current analytics engine evaluates each telemetry data point independently against a static threshold. A \textbf{stateful alerting system} with enter and exit states for each rule would reduce alert noise for slowly varying quantities by triggering alerts only on threshold crossings and introducing hysteresis to prevent rapid toggling.

The edge implementation currently supports a single Raspberry Pi 3B. A centralized \textbf{edge fleet management} system would enable discovery and registration of multiple edge nodes, centralized configuration distribution, health monitoring, and coordinated firmware updates across distributed edge deployments.

The platform currently implements Z-score statistical anomaly detection with configurable sliding windows. Advancing to \textbf{machine learning-based anomaly detection} models, such as isolation forests, autoencoders, or \gls{LSTM} networks, would complement the current Z-score approach by detecting more complex, non-linear, and seasonal patterns that statistical methods may miss. The event-driven architecture supports this naturally: a new ML consumer could subscribe to the \texttt{telemetry.events} Kafka topic alongside existing services without modifying any existing component.

The current edge node evaluates only \texttt{CRITICAL} severity rules to conserve resources. As more capable edge hardware becomes available, this \textbf{expanded edge rule evaluation} could be relaxed to include \texttt{HIGH} or \texttt{MEDIUM} rules, with dynamic severity thresholds based on available hardware resources.

The platform does not currently implement automated \textbf{telemetry retention policies}. A production deployment would require configuring TimescaleDB's built-in retention policies to automatically drop chunks older than a configurable threshold, potentially extended with tiered storage for long-term trend analysis.

The Notifications Service currently implements email delivery with user-configurable severity filtering. Extending the platform with additional \textbf{notification delivery channels}, such as SMS through providers such as Twilio and push notifications to mobile devices, would broaden alert reach. Each additional channel would subscribe to the \texttt{alerts.events} Kafka topic under its own consumer group, maintaining the event-driven decoupling already established by the email delivery implementation.

The current \gls{CoAP} implementation provides basic \gls{UDP}-based telemetry ingestion. \textbf{Advanced \gls{CoAP} features} such as \gls{DTLS} encryption for transport security, the CoAP Observe pattern for server-push telemetry subscriptions, and block-wise transfer for large payloads would strengthen the protocol adapter's capabilities for production constrained-device deployments.

Moving the platform from proof-of-concept to production would require \textbf{production hardening}: \gls{TLS} encryption for all internal communication, proper secret management, Kafka broker clustering with topic replication, TimescaleDB replication with automated failover, and mutual \gls{TLS} between the edge node and the cloud.
